\documentclass[10pt]{article}
\usepackage[margin=0.8in]{geometry}
\usepackage{amsmath,booktabs,hyperref}
\usepackage{xurl}
\usepackage[T1]{fontenc}
\setlength{\parindent}{0pt}
\setlength{\parskip}{5pt}
\title{Why the 0.10-m collision margin permits smaller recorded clearances}
\author{Collision Avoidance Research}
\date{October 3, 2026}
\begin{document}
\maketitle
\textbf{Finding.} The dominant defect is temporal coverage of the collision constraints.
The controller constrains each 50-ms hold at its start and interpolated midpoint,
with the endpoint covered by the next prediction node. The closest configuration
can lie between those nodes. The affine pose trajectory itself crosses below the
requested margin, even when the optimized sampled rows exceed it.

\textbf{Reproduction.} The latest path-bound campaign contains four sub-margin cases.
The diagnostic reproduces all 110 commands through their critical holds exactly,
starting from the recorded measurements. Recorded 31-sample minima agree within
$10^{-9}$ m. Each selected hold is then integrated independently with ODE45
(relative/absolute tolerances $10^{-11}/10^{-12}$), sampled at 1001 times, and
locally refined with scalar minimization. These are numerical diagnostics, not
validated continuous-time bounds. No controller parameter or algorithm changed.

\begin{center}
\begin{tabular}{lrrrr}
\toprule
Speed / scenario & Hold start (s) & Recorded (cm) & Refined (cm) & Affine (cm) \\
\midrule
15 / acceleratingHeadOn & 1.050 & 8.471 & 8.461 & 8.518 \\
8 / crossing & 1.400 & 9.240 & 9.206 & 9.258 \\
8 / curvedCrossing & 1.400 & 9.349 & 9.289 & 9.284 \\
8 / turningCrossing & 1.450 & 5.849 & 5.733 & 5.715 \\
\bottomrule
\end{tabular}
\end{center}
The affine column measures the linear pose interpolant between the optimized
first state and successor; it already falls below 10 cm without nonlinear plant
error. The supplied target prediction retains its exact prescribed motion.

\textbf{Worst case: 8-m/s turning crossing, frame 30.}
At $t=1.450$ s, actual clearance is 11.364 cm; at the constrained midpoint
$t=1.475$ s it is 16.827 cm; at the endpoint $t=1.500$ s it is 38.000 cm.
Nevertheless, at approximately $t=1.460759$ s the refined actual clearance is
5.733 cm. At that same time the affine pose clearance is 5.715 cm. The pose
error, including body rotation, is only 0.316 mm. The first-stage PCBF slack
is $1.24\times10^{-10}$ m. Neither a large slack nor a large first-hold
linearization error explains the roughly 4.27-cm margin shortfall.

\textbf{A smaller, separate approximation defect.}
For the 8-m/s straight crossing at $t=1.425$ s, the linearized row is
9.999967 cm, the affine-pose rectangle gap is 10.000055 cm, and the ODE gap is
9.936264 cm. The one-step RK4 implementation uses an interpolated midpoint;
node/model-geometry errors are not reserved inside the 0.10-m margin. Thus even
an actual midpoint can lie slightly below it. The prediction-agreement function
limits pose/state/CLF error and path deviation, but does not enforce an interval
collision-clearance bound.

\textbf{Design implication.} To interpret the configured margin as a minimum
separation throughout an executed hold, the constraints must cover the hold
interval and account for approximation error. Accurate discrete-time solves
alone do not supply this implication. Closest-time constraints or additional
sampling need an accompanying interval error bound before they can establish
an all-time guarantee. No such correction is implemented or claimed here.

\textbf{Artifacts.} Reproduce with
\path{scripts/auditCollisionMarginSampling.m}; numerical results and source
provenance are under \path{report/MARGIN_SAMPLING_DIAGNOSIS_20261003/}.
The analysis uses simulation commit
\path{8347c2aac7952ae8065f9c571e3c46cdca77e049} and the numerically equivalent
renamed controller at \path{de0767ce8d3bae397c6f2c25d7a32a1a683eaf9e}.
\end{document}
